\documentclass[11pt,a4paper]{article}
\usepackage[utf8]{inputenc}
\usepackage[margin=1in]{geometry}
\usepackage{amsmath,amssymb,amsthm}
\usepackage{listings}
\usepackage{xcolor}
\usepackage{hyperref}

\newtheorem{theorem}{Theorem}
\newtheorem{lemma}[theorem]{Lemma}

\lstset{
  language=C++,
  basicstyle=\ttfamily\small,
  keywordstyle=\color{blue}\bfseries,
  commentstyle=\color{gray},
  stringstyle=\color{red},
  numbers=left,
  numberstyle=\tiny\color{gray},
  breaklines=true,
  frame=single,
  tabsize=4,
  showstringspaces=false
}

\title{IOI 2014 -- Rail}
\author{}
\date{}

\begin{document}
\maketitle

\section{Problem Summary}

There are $n$ railway stations on a line, each of type~C (``left-turn'') or type~D (``right-turn''). Station~0 is type~C at a known position. You may query the distance $d(i,j)$ between any two stations. Using at most $O(n)$ queries, determine the type and position of every station.

The distance function reflects the rail mechanics: trains bounce off intermediate stations of the opposite type, so $d(i,j)$ is not simply $|\text{pos}_i - \text{pos}_j|$ in general.

\section{Solution}

\begin{enumerate}
  \item \textbf{Find the nearest D-station.} Query $d(0, i)$ for all $i \ge 1$. The station $r$ minimising $d(0, r)$ is the nearest D-station to the right of station~0. Its position is $\text{pos}[0] + d(0, r)$.

  \item \textbf{Classify stations.} Query $d(r, i)$ for each remaining station $i$. Using the triangle-inequality test:
  \begin{itemize}
    \item If $d(0, i) = d(0, r) + d(r, i)$: station $i$ lies to the right of $r$ (candidate C-station).
    \item Otherwise: station $i$ lies to the left of station~0 (candidate D-station).
  \end{itemize}

  \item \textbf{Verify and determine positions.} Process right-side candidates sorted by $d(r, i)$ (ascending), tracking the last confirmed D-station. For each candidate, one verification query against the last confirmed D-station determines its type and position. Symmetrically, process left-side candidates sorted by $d(0, i)$, tracking the last confirmed C-station.
\end{enumerate}

\section{C++ Implementation}

\begin{lstlisting}
#include <bits/stdc++.h>
using namespace std;

// Grader-provided: int getDistance(int i, int j);

void findLocation(int n, int first, int location[], int stype[]) {
    stype[0] = 1;
    location[0] = first;
    if (n == 1) return;

    vector<int> d0(n);
    d0[0] = 0;
    for (int i = 1; i < n; i++)
        d0[i] = getDistance(0, i);

    // Nearest station to 0 is the closest D-station to its right
    int r = 1;
    for (int i = 2; i < n; i++)
        if (d0[i] < d0[r]) r = i;
    stype[r] = 2;
    location[r] = first + d0[r];

    vector<int> dr(n, -1);
    vector<int> rightC, leftD;
    for (int i = 1; i < n; i++) {
        if (i == r) continue;
        dr[i] = getDistance(r, i);
        if (d0[i] == d0[r] + dr[i])
            rightC.push_back(i);
        else
            leftD.push_back(i);
    }

    // Process right-side C-candidates
    sort(rightC.begin(), rightC.end(),
         [&](int a, int b) { return dr[a] < dr[b]; });
    int lastD = r;
    for (int i : rightC) {
        int expected_pos = location[r] + dr[i];
        int dist_check = getDistance(lastD, i);
        int direct = expected_pos - location[lastD];
        if (dist_check == direct && direct >= 0) {
            stype[i] = 1;
            location[i] = expected_pos;
        } else {
            stype[i] = 2;
            location[i] = location[lastD] - dist_check;
            lastD = i;
        }
    }

    // Process left-side D-candidates
    sort(leftD.begin(), leftD.end(),
         [&](int a, int b) { return d0[a] < d0[b]; });
    int lastC = 0;
    for (int i : leftD) {
        int expected_pos = location[0] - d0[i];
        int dist_check = getDistance(lastC, i);
        int direct = location[lastC] - expected_pos;
        if (dist_check == direct && direct >= 0) {
            stype[i] = 2;
            location[i] = expected_pos;
        } else {
            stype[i] = 1;
            location[i] = location[lastC] + dist_check;
            lastC = i;
        }
    }
}
\end{lstlisting}

\section{Complexity Analysis}

\begin{itemize}
  \item \textbf{Queries}: $(n-1)$ for $d(0,\cdot)$ + $(n-2)$ for $d(r,\cdot)$ + at most $(n-2)$ verification queries $= 3n - 5$ total, i.e., $O(n)$.
  \item \textbf{Time}: $O(n \log n)$ due to sorting.
  \item \textbf{Space}: $O(n)$.
\end{itemize}

\end{document}
